% Slajd 9 – model biznesowy i komercjalizacja (kryterium: 20 %; pkt 9 wyzwania).
\begin{frame}{Model biznesowy i~komercjalizacja}
  \footnotesize
  \renewcommand{\arraystretch}{1.1}
  \begin{tabularx}{\textwidth}{@{}>{\bfseries\raggedright\arraybackslash}p{3.1cm} X >{\raggedright\arraybackslash}p{3.4cm}@{}}
    \toprule
    Segment & Oferta & Przychód \\
    \midrule
    Mieszkańcy i~turyści & aplikacja, profil potrzeb, asystent, mapa barier & bezpłatnie, zawsze -- to oni tworzą dane \\
    Właściciele obiektów (hotele, gastronomia, kultura, sieci) & panel podstawowy bezpłatny; \textbf{plan Pro}:
      odznaka właściciela, statystyki „czego szukają goście, czego brakuje”, hurtowe udogodnienia,
      widget dostępności na stronę i~do systemu rezerwacji & abonament za obiekt lub pakiet dla sieci \\
    Organizatorzy wydarzeń, zarządcy nieruchomości & pakiet czasowy: tymczasowe wejścia i~trasy, mapa barier
      na czas imprezy & opłata za wydarzenie / obiekt \\
    Miasta i~regiony & wdrożenie white-label (konfiguracja, import OSM, warstwy miejskie); raporty
      o~brakach danych i~udogodnień & wdrożenie + roczne utrzymanie \\
    Mapy, aplikacje turystyczne, systemy rezerwacyjne & Open API z~atrybucją: fakty o~dostępności
      ze źródłem, datą i~statusem & licencja API / limity zapytań \\
    \bottomrule
  \end{tabularx}

  \vspace{1mm}
  \textbf{Zasady, które chronią wartość danych.} Nie sprzedajemy danych osobowych; treści społeczności i~dane OSM
  pozostają otwarte (ODbL). Właściciel płaci za narzędzia i~widoczność, nie za „lepszy” status -- status zawsze wynika
  z~weryfikacji.

  \note[item]{Najbardziej realny pierwszy przychód: hotele i~turystyka medyczna (profil „Po zabiegu”, warstwa SOR / apteki, 4~języki) – goście zza granicy pytają o~dostępność przed rezerwacją.}
  \note[item]{Panel właściciela już dziś pokazuje „Najczęściej wyszukiwane informacje o~Twoim lokalu” – to zalążek statystyk z~planu Pro.}
  \note[item]{Pakiet dla wydarzeń: w~panelu właściciela są już punkty w~obiekcie (windy, toalety, wejścia) dodawane hurtowo – organizator wnosi je na czas imprezy.}
  \note[item]{Miasto nie jest klientem w~pierwszej kolejności; korzysta z~danych i~raportów, a~ewentualnie finansuje wdrożenie white-label.}
\end{frame}
